\section{Strategy of the proof}\label{S:strategy}

The second moment is a convergent sum of two-chain pairing integrals,
including empty chains and all relative time orders. Stratonovich integrals
retain pairings within either chain as well as between the chains, so products
of different orders can contribute (Proposition~\ref{P:mixed}). Estimating
every pairing by one common bound gives a finite moment but loses the target
growth scale (Proposition~\ref{P:envelope}). Chapter~\ref{S:book-no-go}
discusses this and the other approaches that do not work. The proof proceeds in
four steps, shown in Figure~\ref{Fig:proof-map}.

\begin{figure}
  \centering
  \begin{tikzpicture}[x=1cm,y=1cm,
    destination/.style={ham/box,text width=5.65cm,minimum height=1.55cm},
    result/.style={destination,draw=hamteal,fill=hamteal!6}]
    \node[destination] (construction) at (0,0)
    {\textbf{The canonical solution}\\
    \GuideRef{Chapter}{S:G0}\\
    \GuideRef{Proposition}{P:mild-construction}};
    \node[destination] (upper) at (-3.35,-2.4)
    {\textbf{Independent upper bound}\\
    \GuideRef{Chapter}{S:book-independent-upper}\\
    \GuideRef{Theorem}{T:full-upper}};
    \node[destination] (phase) at (3.35,-2.4)
    {\textbf{Noise and full wave phase}\\
    \GuideRef{Ch.}{S:book-gaussian-foundations},
    \GuideRef{Ch.}{S:book-operator-paths},
    \GuideRef{Ch.}{S:book-wave-evolution}\\
    \GuideRef{Section}{S:constant-phase-young}};
    \node[result] (finite) at (-3.35,-5)
    {\textbf{Finite variational value}\\
    \GuideRef{Section}{S:unmasked-variational}\\
    \GuideRef{Proposition}{P:variational-positive-scores}};
    \node[destination] (lower) at (3.35,-5)
    {\textbf{Lower bound for each control}\\
    \GuideRef{Chapter}{S:book-lower}\\
    \GuideRef{Proposition}{P:raw-block-original-lower}};
    \node[destination] (signed) at (0,-7.6)
    {\textbf{The signed-control upper bound}\\
    \GuideRef{Ch.}{S:book-means},
    \GuideRef{Ch.}{S:book-growing-radius},
    \GuideRef{Ch.}{S:book-point-upper}\\
    \GuideRef{Proposition}{P:original-signed-point-upper}};
    \node[destination] (comparison) at (0,-10.1)
    {\textbf{Positive Gram comparison}\\
    \textbf{and cutoff removal}\\
    \GuideRef{Section}{S:finite-positive-gram}\\
    \GuideRef{Proposition}{P:positive-score-cutoff}};
    \node[result] (limit) at (0,-12.6)
    {\textbf{Matching upper and lower values}\\
    $C_2(\Gamma,\alpha_0)=S_2^{\mathrm{pos}}$\\
    \GuideRef{Chapter}{S:book-constant};
    \GuideRef{Theorem}{T:intrinsic-second-moment}};
    \draw[ham/arrow] (construction)--(upper);
    \draw[ham/arrow] (construction)--(phase);
    \draw[ham/arrow] (upper)--(finite);
    \draw[ham/arrow] (phase)--(lower);
    \draw[ham/arrow] (lower)--(finite);
    \draw[ham/arrow] (finite)--(signed);
    \draw[ham/arrow] (signed)--(comparison);
    \draw[ham/arrow] (comparison)--(limit);
    \draw[ham/arrow] (lower.east)--++(.75,0) |- (limit.east);
  \end{tikzpicture}
  \caption[Proof dependencies and destinations]{The principal proof
    dependencies, with links to chapters and results in the technical companion.
    The independent
    upper bound and the fixed-control lower bound first give a finite
    variational lower value. Further operator and Gaussian estimates give a
    signed-control upper bound for the same point moment. Positive Gram
    comparison and cutoff removal identify the two values. Arrows show
    principal uses; each box also has the hypotheses and earlier inputs stated
    at its linked destination. They do not assert independence of successive
    time blocks. The zero-covariance case is included in the final theorem.}
  \label{Fig:proof-map}
\end{figure}


\paragraph{One solution and an upper bound at the right scale.}
Chapter~\ref{S:G0} derives the complete pairing expansion and constructs the
solution used throughout (Proposition~\ref{P:mild-construction}).
Chapter~\ref{S:book-independent-upper} then proves an upper bound at the scale
$t^\kappa$ for the full covariance class, before the limiting constant is known
(Theorem~\ref{T:full-upper}). The wave equation forces this argument to carry
both displacement and velocity across time intervals. The constant of this
bound is not claimed to be optimal.

\paragraph{A lower bound for each control.}
A lower bound comes from shifting the Gaussian noise toward a deterministic
potential that makes the wave grow, and comparing the gain with the cost of the
shift. Because the noise is correlated in time, the cost of a control is
charged in the original covariance space, not added over time intervals as if
they were independent. Chapters~\ref{S:book-gaussian-foundations}--\ref{S:book-wave-evolution}
prepare the tools: covariance estimates for a potential multiplied by a smooth spatial cutoff of compact support, a
random operator path driven by the noise of the original equation, and the
evolution on the displacement--velocity phase space, identified with the
physical wave. Chapter~\ref{S:book-lower} converts long time into an
equivalent small-noise problem by equality in law
(Section~\ref{S:exact-canonical-scaling}), for nonzero covariance, proves a lower bound for each fixed control in the
class defined in Section~\ref{S:targets}, with an eligible fixed spatial mask
(Proposition~\ref{P:raw-block-original-lower}), and then removes that mask and
optimizes over controls to obtain a variational lower value
(Section~\ref{S:unmasked-variational}). For zero covariance, the solution is
identically one and the rate is zero. Together with the upper bound of the
first step, this value is finite (Proposition~\ref{P:variational-positive-scores}).

\paragraph{A matching upper bound.}
The upper bound must account for all Gaussian fluctuations, not only the
favourable shifts, and it must control the moment at one point while the region
reached by the wave grows with time. Chapter~\ref{S:book-means} relates limits of smooth controls in the norm determined by the noise covariance to the deterministic means they induce. Chapter~\ref{S:book-growing-radius} separates the estimates that are
uniform in the radius of the spatial cutoff from those whose dependence on the radius must be shown, and constructs the wave for completed static mean potentials of either sign at a fixed spatial cutoff. Chapter~\ref{S:book-point-upper} reduces the probabilistic estimate to
growth under deterministic controls, returns to the original point moment through finite propagation and a spatial average of squared solutions, controls the Gaussian tail, and obtains an upper
variational expression over signed controls
(Proposition~\ref{P:original-signed-point-upper}).

\paragraph{Identifying the constant.}
The two bounds now concern the same point moment, but over different controls:
nonnegative smooth controls below, signed controls in the covariance-norm completion just described above.
Chapter~\ref{S:book-constant} shows that the two optimizations have the same
value. It first solves a finite optimization problem whose matrix of pairwise covariance inner products has nonnegative entries (Section~\ref{S:finite-positive-gram}), then applies it to
the physical means and removes the spatial cutoff
(Proposition~\ref{P:positive-score-cutoff}). Theorem~\ref{T:intrinsic-second-moment}
states the resulting limit. The remaining sections give bounds for its constant
and, for spatial white noise in one dimension, explicit moment bounds; they do
not evaluate the full supremum.
